
\begin{proposition}[Application to the named six-sphere families]\label{cdp:st:E7}
Every member of the following torus-fibred families,
including the period family $X_c$, its members with $|p|>1$,
Viaclovsky's $X_{3k,4l}$, his two-parameter family and his combs,
and the Liu--Rollenske family, satisfies the rank hypothesis of \ref{cdp:st:E6}.
Consequently every sufficiently small deformation retains a
holomorphic torus fibration over $\mathbb P^1$. Every such member
with $b_2=0$, and every sufficiently small deformation of it, has
algebraic dimension one. In particular the fibration-retention
question in Section~10(1) of the supplied comb paper has an
affirmative answer.
\end{proposition}
\begin{proof}
We verify a local condition common to these actual constructions.
At their hexagonal fibre, a marking on the punctured disc puts
the Hodge periods in the form
\begin{equation}\label{cdp:eq:E-log-period}
  (I_2,Z(q)),\qquad
 Z(q)=\frac{\log q}{2\pi i}B+H(q),
 \qquad \det B\ne0,
\end{equation}
where $H$ is holomorphic at zero. Therefore
\[
 qZ'(q)=\frac{B}{2\pi i}+qH'(q)
\]
is invertible for all sufficiently small nonzero $q$. The
Kodaira--Spencer tensor of a marked torus with periods $(I,Z)$
is represented by $Z'$ after the invertible identifications of
the Lie algebra and $H^1(\mathcal O)$ with the chosen period
coordinates. Thus it has rank two at those regular fibres.

Here the existence and nonsingularity of $B$ are checked directly
in the supplied local models, including
\cite[Proposition labelled \texttt{hexfill}]{Viaclovsky2026Combs},
\cite[Proposition labelled \texttt{mum}]{Viaclovsky2026Two}, and
\cite[Proposition labelled \texttt{prop: compactification X}]{LiuRollenske2026}.
The AC cusp has logarithmic matrix
$B_0=\left(\begin{smallmatrix}0&1\\-1&0\end{smallmatrix}\right)$.
The affine twists and finite logarithmic transformations change
the finite collars and constant gluing data, while the punctured
cusp matrix still has this same invertible logarithmic coefficient.
Engel's two deck transformations at the $I_1$-fibre have independent
valuation vectors $(1,0)$ and $(0,1)$, again giving an invertible
$B$. In the two-parameter family the designated fibre over $p_3$
is the full-lattice rank-two filling; its displayed two valuation
vectors are independent. Liu--Rollenske's compactification over
$I_1$ uses the same two independent valuation directions. Each
comparison is a change of the integral marking and of the
holomorphic fibre frame, so preserves the rank of the
Kodaira--Spencer tensor.

For the combs, \cite[Proposition labelled \texttt{periodextension},
case (i)]{Viaclovsky2026Combs} writes a flat frame as
$e_i=\alpha_i+\ell(q)\beta_i$, $f_i=\beta_i$, with
$\ell(q)=\log q/(2\pi i)$, and the extended Hodge subbundle as
the rows of $(I,H(q))$ in this frame. The fixed value of the
smoothed period section at $q=0$ keeps it transverse to the
$f_i$-plane. Its rows therefore remain $(I,H_s(q))$ with
$H_s$ holomorphic. In the flat frame they are
$(I,\ell(q)I+H_s(q))$, which is \eqref{cdp:eq:E-log-period} with
$B=I$. This verifies the needed condition for every comb member,
without a formula for its global splitting type.

All the sources construct the indicated members as compact
connected-fibre holomorphic maps with smooth general two-torus
fibre. \ref{cdp:st:E6} applies, and \ref{cdp:st:E1} gives the algebraic-dimension assertion.
On the differentiable six-sphere these neighbourhoods of the
listed structures constitute an open set on which $a=1$.
\end{proof}


